\section{Discussion}

The two tasks point to the same practical distinction. Calendar patterns explain when rain is more common across many years, whereas recent observations help distinguish the conditions surrounding a particular prediction time. The afternoon peak and seasonal variation therefore justify calendar features, but they do not imply that a calendar model can identify an individual storm. B's lack of overall improvement over historical frequency makes that limit visible. Model A has more useful short-term information, although its advantage narrows for large accumulations.

Missingness affects both conclusions. For exploration, coverage rules leave 85 of 96 year-months with at least one eligible station, while full-year summaries are supported in only three years. For prediction, A's history requirement retains 89.9\% of B's eligible 2024 test cases. Performance on that subset does not establish what would happen during excluded periods. Neither a complete-window rule nor the modest sensitivity of monthly means to coverage cutoffs makes missing observations random.

There are also limits to reconstructing a forecast retrospectively. We treat the observation timestamp as the end of its five-minute accumulation interval; this is an alignment assumption that should be checked against an authoritative measurement convention. Requiring both update timestamps to be no later than the forecast time gives a conservative availability rule, but the annual files do not provide a complete archive of past revisions. For B, the saved evaluation uses coordinates attached to each historical observation, while the prediction interface uses a catalogue frozen from training. A station relocation could make those two representations differ. All 70 evaluated station IDs occur in the training catalogue, so performance at unseen stations is also untested. These issues matter when transferring the offline experiment into an operational service.

A learning curve would show whether the 15\% training sample captures most of the available predictive benefit. A second experiment, fitting A and B to exactly the same station-time examples, would separate the contribution of rainfall history from differences in training populations. Thresholds inside the amount bins should also be evaluated: monotone probabilities are desirable, but monotonicity does not verify the accuracy of the global within-bin approximation.

The present uncertainty estimates are conditional on fitted models and one test year. Rolling-year evaluation would show whether the findings persist across different weather and station conditions; changing the time-block length would test the sensitivity of the reported interval. None of these extensions has been completed here. We therefore interpret the results as evidence for the observed network in 2024, with the strongest support for lower-threshold short-term prediction when recent station observations are available.

\section{Conclusion}
The records show clear seasonal, diurnal and local variation, but their gaps prevent a uniform account of annual rainfall across all eight years. Coverage-aware monthly summaries and common-station comparisons preserve useful evidence without treating missing periods as dry weather. For prediction, the recent-rainfall forest improves substantially on the calendar forest and the within-mode baselines in mean probability error. The remaining difficulty is most evident for rare, large accumulations. A usable probability estimate therefore needs both an appropriate model and a clear statement of the observations, threshold range and population to which its evaluation applies.

\section*{AI usage declaration}
OpenAI Codex assisted with code development, data-cleaning diagnostics, cloud execution, feature and model implementation, tests, figure preparation, numerical analysis and report drafting. It also assisted with restructuring and editing this report and redrawing the calibration and December station figures from saved results. Reported numerical results come from the recorded data analyses and model outputs; synthetic records were used only for software tests. Group members remain responsible for checking the code, figures, interpretations and references in the submitted work.
